\section{Related Work}

The evolution of digital advertising has transitioned from static heuristic-based bidding to dynamic, data-driven optimization. However, existing paradigms often struggle with the dual challenge of exploration-exploitation trade-offs and the inherent non-stationarity of user behavior. In this section, we contextualize \textit{Active-Bidder} within the landscape of Click-Through Rate (CTR) prediction, Reinforcement Learning (RL) for bidding, and the emerging field of Active Inference.

\subsection{CTR Prediction and Real-Time Bidding (RTB)}
Early approaches to RTB relied heavily on pointwise estimation of Click-Through Rates using Logistic Regression or Gradient Boosted Decision Trees (GBDT). With the advent of deep learning, models like DeepFM \citep{guo2017deepfm} and Wide \& Deep \citep{cheng2016wide} sought to capture high-order feature interactions. The introduction of the Transformer architecture \citep{vaswani2017attention} catalyzed a shift toward sequential modeling in recommendation systems, enabling models to capture long-term user interests via self-attention mechanisms \citep{zhou2018deep}.

In the current era, the field is moving toward Large Recommendation Models (LRMs) that leverage scaling laws similar to those observed in Large Language Models (LLMs) \citep{dubey2024llama, zhao2023survey}. Recent works have explored the application of Foundation Models for generalist decision-making \citep{wang2024generalist, pathak2024foundation}, yet these models typically optimize for a surrogate LogLoss objective. This objective, while effective for ranking, fails to account for the \textit{market dynamics} of the auction, where the optimal bid is a function of both the predicted value and the expected competition. Unlike these point-wise or sequence-to-sequence estimators, \textit{Active-Bidder} treats the bidding process as a generative task, explicitly modeling the uncertainty of the market environment.

\subsection{Reinforcement Learning in Advertising}
To move beyond static prediction, Reinforcement Learning has been widely adopted to optimize bidding strategies in multi-step environments \citep{sutton2018reinforcement}. Deep Reinforcement Bidding (DRB) and its variants have successfully formulated RTB as a Markov Decision Process (MDP), using architectures like Deep Q-Networks (DQN) or Soft Actor-Critic (SAC) to maximize long-term Total Value or Return on Ad Spend (ROAS) \citep{wu2018deep}.

However, a fundamental limitation in current RL-based advertising research is the \textbf{Passive User Assumption}. Most RL agents treat the user and the ad auction as a black-box environment that merely provides rewards (clicks/conversions) in response to actions (bids) \citep{silver2016mastering}. This ignores the fact that users are active agents whose internal states (e.g., intent, fatigue) evolve in response to exposure. Furthermore, standard RL often relies on naive exploration strategies like $\epsilon$-greedy or entropy regularization, which are inefficient in the high-dimensional, sparse-reward settings of RTB. Recent advances in offline RL and conservative Q-learning \citep{levine2020offline} attempt to mitigate training instability, but they do not address the epistemic gap—the uncertainty about the environment's underlying causal structure. \textit{Active-Bidder} diverges from these approaches by adopting a world-model-based perspective, similar to recent developments in Test-Time Training (TTT) and adaptive game theory \citep{sun2024learning, yang2024zero}.

\subsection{Active Inference and Generative Control}
Active Inference, rooted in the Free Energy Principle (FEP) \citep{friston2010free}, offers a mathematically rigorous framework for unifying perception and action. Unlike RL, which maximizes an external reward signal, Active Inference agents minimize \textit{variational free energy}, which naturally decomposes into extrinsic value (reward seeking) and epistemic value (information seeking/exploration).

While Active Inference has traditionally been studied in the context of cognitive science and robotics, its application to socio-technical systems like digital advertising is nascent. Recent work in ``generative control'' has shown that agents equipped with a generative model of their environment can outperform model-free agents in non-stationary tasks \citep{peebles2023scalable}. The integration of geometric priors \citep{bronstein2024geometric} and neural operators \citep{kovachki2023neural, lu2024physics} has further enhanced the ability of such models to handle complex, continuous state spaces.

The core contribution of \textit{Active-Bidder} lies in its formulation of the bidding agent as an ``Active Inferer'' that maintains a belief over the hidden state of the user. This allows for \textit{directed exploration}—the agent may choose to bid on a sub-optimal impression specifically because it yields the most information about the user’s current preference, thereby reducing future uncertainty and maximizing long-term ``surprisal'' minimization. This aligns with recent trends in LLM-based agents that utilize test-time adaptation to handle variable-length contexts and evolving environments \citep{reid2024llama, sun2024learning}.

\subsection{Summary of Differentiation}
Compared to SOTA baselines, \textit{Active-Bidder} introduces three critical innovations:
\begin{enumerate}
    \item \textbf{Beyond LogLoss}: We replace simple point-wise prediction with a variational generative model that captures the joint distribution of market price and user response.
    \item \textbf{Active User Modeling}: We dismantle the ``Passive User'' assumption by modeling the user as a dynamical system whose latent state is inferred through interaction, rather than a static vector.
    \item \textbf{Intrinsic Motivation}: We provide a formal bridge between the FEP and RTB, using the Expected Free Energy (EFE) as a unified objective for bidding, which naturally balances budget constraints with the need for market exploration.
\end{enumerate}
By leveraging recent breakthroughs in scalable transformers \citep{dao2023flashattention} and selective state-space models \citep{gu2023mamba}, \textit{Active-Bidder} achieves a level of adaptive efficiency previously unattainable in large-scale bidding systems.